% 06-massless-limit — the analytic anchor; kept short on purpose
\section{Massless mesons and baryons in the $m/g \to 0$ limit}
\label{sec:massless}

At $m = 0$ the spectrum contains towers of exactly massless hadrons.  The
construction is Sec.~VI of Ref.~\cite{Hornbostel:1988fb} and survives
unchanged, so we record only what this paper uses.  The color currents
$V^a_k$ obey a Kac--Moody algebra, and the interacting part of $P^-$ is a
quadratic form in them.  States created purely by currents, and for baryon
number by a composite baryon operator, are annihilated by it.  Massless
mesons, baryons, and multi-hadron states therefore exist at every
resolution $K$, not merely in the limit.  The massless baryon's valence
distribution is closed-form,
\begin{equation}
  q(x) \;=\; N(N-1)\,(1-x)^{N-2},
  \label{eq:massless-sf}
\end{equation}
and collapses at $N = 2$ to the flat meson distribution.  For U($N$) the
extra U(1) current promotes these states to free bosons at
$M^2 = g^2/2\pi$: the Schwinger mechanism, and the sharpest statement of
how SU($N$) and U($N$) differ in two dimensions.

Three later results lean on this section.  The identically-zero last row of
Table~\ref{tab:tableone} is exact and $K$-independent; our solvers
reproduce it at every resolution as a structural check.  The chiral fits of
Sec.~\ref{sec:chiral} use masslessness as the fixed endpoint any
$M^2(m/g)$ form must respect.  And the exotics search of
Sec.~\ref{sec:exotics} inherits a boundary condition: multi-hadron states
are exactly massless at $m = 0$, so every binding energy we measure must
vanish in the chiral limit.  A candidate whose $\Delta$ fails to do so is
diagnosing its own systematics, not new physics.
